\section{Einleitung}

\subsection{Motivation und Problemstellung}

Die Rekonstruktion dünner, linearer Objekte aus Bildsequenzen verbindet mehrere
schwierige Teilprobleme: Die Kamerabewegung muss stabil geschätzt werden, das
Objekt muss über die Bilder hinweg konsistent segmentiert werden und die
objektbezogene 3D-Repräsentation muss für nachgelagerte Geometrieoperationen
verwendbar sein. Die Projektarbeit untersucht diesen Ablauf zunächst an einem
kontrollierten Alurohr-Testdatensatz. Die spätere Bachelorarbeit soll die
Pipeline auf reale Rohr- oder Kabeltrassen und geodätische Referenzen
übertragen.

\subsection{Zielsetzung und Eigenbeitrag}

Das Ziel der Projektarbeit ist ein technischer Machbarkeitsnachweis. Dazu muss
die Pipeline vom Video bis zur lokalen Centerline durchlaufen, die zentralen
Datenverträge müssen überprüfbar sein und die wichtigsten Varianten müssen mit
reproduzierbaren Parametern verglichen werden. Eine Aussage über die Einhaltung
einer Toleranz von \(\pm 10\,\mathrm{cm}\) wird für die Projektarbeit nicht
behauptet, weil dafür noch geeignete GNSS- oder Vermessungsreferenzen fehlen.
Der Eigenbeitrag liegt in der Integration der heterogenen Komponenten, der
Trennung von Roh- und Idealbilddomäne (die verzeichnungsbehafteten Rohbilder
werden in verzeichnungsfreie PINHOLE-Bilder überführt, sodass Masken,
Kamerageometrie, Splatting und Auswertung in derselben, konsistenten Domäne
arbeiten), der Definition überprüfbarer Datenverträge (Schnittstellenregeln,
die festlegen, welche Artefakte in welcher Bilddomäne zueinander passen
müssen und nach welchen Kriterien eine Stufe als gültig gilt), der
maskengetriebenen Objekt- und Meshgewinnung sowie der automatisierten
Durchführung und wissenschaftlichen Interpretation der Tests. STS und SuGaR
sind dabei genutzte, nicht selbst entworfene Verfahren; der eigene Anteil in
der Segmentierung liegt in der angepassten Maskenlogik: die Erosionshierarchie
(default unverändert, middle und small erodiert) sowie die promptbasierte
Auswahl sorgen dafür, dass automatisch genau ein
segmentiertes Zielobjekt entsteht. Das offizielle STS-Verfahren geht davon
verschieden, indem es sämtliche Szenenobjekte automatisch in drei
Granularitätsstufen erfasst und erst nach der Rekonstruktion eine offene
Textabfrage ermöglicht; der Objektprompt legt das Ziel hier dagegen bereits
vor dem Training fest, sodass Mehrfachsegmentierung und Nachabfrage
entfallen. Die Funktionsweise ist in
Abschnitt~\ref{sec:grundlagen} beschrieben.
Der Einsatz von SAM~3.1 verdoppelt zudem die Verarbeitungsgeschwindigkeit für
Videos mit mittlerer Objektanzahl (16 auf 32 Frames pro Sekunde auf einer
einzelnen H100-GPU) \parencite{metaSam3}.

\subsection{Forschungsfragen}

Die zentrale Arbeitsfrage lautet:
\begin{quote}
Lässt sich eine Docker-orchestrierte, objektbezogene Scan-to-BIM-Pipeline aus
promptbarer SAM-3.1-Segmentierung, kameramodellbewusstem COLMAP-SfM,
Segment-then-Splat und meshbasierter Nachverarbeitung für ein lineares
Testobjekt robust und reproduzierbar bis zu einer lokalen Centerline betreiben?
\end{quote}

Zur Beantwortung werden vier Unterfragen betrachtet:

\begin{enumerate}
  \item Welche Bilddomänen, Datenverträge und Quality-Gates sind erforderlich,
	  damit Masken, Kameras, Eval-Split, Gaussians, Mesh und Centerline
	  konsistent übergeben werden?
  \item Wie beeinflussen Kameramodell, zeitliche Abtastung und Bildauflösung
	  den technischen Pipelineerfolg und die objektmaskierten Bildmetriken?
  \item Welche der untersuchten Meshrouten -- direkte Original-STS-GS-Route A
	  oder SuGaR-Coarse -- bewahrt die für die Centerline benötigte
	  Objektgeometrie im Projektstand geeigneter?
  \item Inwieweit ermöglichen interaktiver CLI-Betrieb, Replay, Autopilot und
	  Matrixrunner reproduzierbare Einzel- und Batchläufe?
\end{enumerate}

\subsection{Abgrenzung und Aufbau}

Die Arbeit konzentriert sich auf Grundlagen, Konzept, Implementierung,
Funktions- und Vergleichstests sowie die Interpretation der Ergebnisse. Lange
Installationsanleitungen, vollständige Quellcode-Dokumentation und spätere
Bachelorarbeitsfragen werden bewusst außerhalb des Haupttextes gehalten. Nicht
Gegenstand sind eine reale Vermessung einer Kabeltrasse, ein Nachweis von
$\pm10\,\mathrm{cm}$, eine vollständige Bewertung aller Kameramodelle, eine
allgemeine Einführung in Computer Vision oder eine Produktionsentscheidung
allein aus PSNR, SSIM und LPIPS. Die vorhandene Web-UI bleibt ebenfalls
außerhalb des Untersuchungsumfangs und ist der späteren Bachelorarbeit
vorbehalten; betrachtet werden CLI, Autopilot, Replay
und Matrixrunner.
